\section*{Chapter \ref{ch:in:structurer-sales-and-sales-trader} --- Structurer, Sales and Sales-Trader}

\begin{solution}{exo:in:structurer-sales-and-sales-trader:1}
$(286\,264+99\,890)/486\,635=79.4\%$.
\end{solution}

\begin{solution}{exo:in:structurer-sales-and-sales-trader:2}
A margin of 1.83 per 100; a coupon of 6.97\% a year would leave none.
\end{solution}

\begin{solution}{exo:in:structurer-sales-and-sales-trader:3}
Article 16(3): a process for approving each financial instrument before it is marketed, specifying an identified target
market of end clients, assessing the risks to it and checking the distribution strategy against it, with regular review.
\end{solution}

\begin{solution}{exo:in:structurer-sales-and-sales-trader:4}
The note is a bond plus options the investor sells. At a higher rate the bond part costs less today for the same
repayment, which frees value for the coupon; the forward of the index also rises, which makes the autocall more likely.
\end{solution}

\begin{solution}{exo:in:structurer-sales-and-sales-trader:5}
He sells more volatility: at 30\% the chance that the index ends below 60\% and he loses capital rises from 9.4\% to
19.3\% of paths, and the coupon pays for that.
\end{solution}

\begin{solution}{exo:in:structurer-sales-and-sales-trader:6}
The desk books the margin when the note is sold and hedges the note's risk from then on; what the client later earns or
loses depends on the market, which the \omterm{def:in:structurer-sales-and-sales-trader:structurer}{structurer} does not control. The link to the client's outcome runs instead
through \omterm{def:in:structurer-sales-and-sales-trader:governance}{product governance}, the desk's reputation and future sales.
\end{solution}

\begin{solution}{exo:in:structurer-sales-and-sales-trader:7}
Across the 20 seeds the margin averages 1.84 per 100 with a standard deviation of 0.06 (range 1.74 to 1.97): its sign is
not in doubt, but the second decimal is.
\end{solution}

\begin{solution}{exo:in:structurer-sales-and-sales-trader:8}
The protection holds only if the index ends above 60\%; below it the holder bears the whole fall from the start, so
the note has the index's worst outcomes without its best (it is called away when the index rises). The coupon is not
guaranteed either (it needs the index above 70\%), and the holder forgoes the dividends. Compare the loss probability and
the distribution, not the headline coupon.
\end{solution}

\begin{solution}{pb:in:structurer-sales-and-sales-trader:1}
\begin{enumerate}
\item As in the chapter's definitions.
\item The client or distributor, the trading desk that hedges the note, and the bank's margin.
\item Works clients' orders in the market and reports against their benchmark; order handling with a person in the loop
(Book~10, chapter~20).
\item Pricing and sensitivities; translation of views and products; execution measurement.
\item Book~5, chapter~18.
\item Five years, annual observations, autocall at 100\%, 6\% coupon at 70\% with memory, capital lost below 60\% at
maturity.
\item $100-V(c)$; $c_m=(100-m-V(0))/(\partial V/\partial c)$ because the value is linear in the coupon.
\item $-1.65$, 1.83, 6.91 and 9.13 per 100.
\item 4.23\%, 5.91\%, 8.79\% and 10.22\%.
\item The investor sells the downside and is paid for it; the probability of losing capital rises too.
\item 486\,635 million euros in six markets; Switzerland 58.8\%, Germany 20.5\%.
\item 2 to 10 million euros a year at a 2\% margin on 1 to 5 billion euros of issuance and ten \omterm{def:in:structurer-sales-and-sales-trader:structurer}{structurers}.
\item Articles 16(3) and 24(2) (approval and target market) and 25(1) (knowledge and competence), with ESMA's 2023
guidelines.
\item The Series 7 qualifies a general securities representative, with the SIE; the Series 63 covers state securities law.
\item Sales agents' median \$103\,030 in the securities industry (90th percentile \$309\,440), \$61\,440 at banks.
\item At 6\%: $-1.65$ at 1\% and 1.83 at 4\% (20\% volatility); 6.91 and 9.13 at 30\%; 2 to 10 million euros per
\omterm{def:in:structurer-sales-and-sales-trader:structurer}{structurer}.
\item It raises the margin at a given coupon, or the coupon at a given margin, and is also paid by the buyer.
\item A product that fits the target market, a description of its risks the client can understand, and a fair price.
\item Pricing, risk and sensitivity to parameters the client does not see; not the relationships and distribution that
decide volume.
\item The desk's revenue is margin times volume, and the margin depends on rates and volatility as much as on design.
Take the job if you want to work with clients and products; your models will matter most in saying what a note really
sells.
\end{enumerate}
\end{solution}

\begin{solution}{iq:in:structurer-sales-and-sales-trader:1}
A zero-coupon bond paying 100 at maturity plus a call option on the index (participation), bought with the difference
between 100 and the bond's price; the issuer's margin comes out of that difference.
\iqlookfor{bond plus option and the budget.}
\end{solution}

\begin{solution}{iq:in:structurer-sales-and-sales-trader:2}
Higher rates (a cheaper bond part), higher implied volatility (the downside sold is worth more), and a different design
(a lower barrier or a worst-of); the margin may also differ.
\iqlookfor{rates, volatility and design.}
\end{solution}

\begin{solution}{iq:in:structurer-sales-and-sales-trader:3}
Short volatility (short the downside put); in a worst-of, short correlation: lower correlation makes the worst performer
worse and the note cheaper, so the coupon is higher.
\iqlookfor{the sign of each exposure and why.}
\end{solution}

\begin{solution}{iq:in:structurer-sales-and-sales-trader:4}
Compare against the market's move during the order and the expected impact for its size: adverse drift, spread paid,
impact of 5\% participation; then show whether the algorithm behaved as designed, with the transaction cost analysis.
\iqlookfor{decomposing slippage into market move, spread and impact.}
\end{solution}

\begin{solution}{iq:in:structurer-sales-and-sales-trader:5}
Client type, knowledge and experience, ability to bear loss (up to the whole capital below 60\%), risk tolerance,
objectives (income with conditional protection), horizon (up to five years), and the negative target market.
\iqlookfor{the MiFID target-market categories applied to the note.}
\end{solution}

\begin{solution}{iq:in:structurer-sales-and-sales-trader:6}
The desk is long volatility near the barriers and long the downside the investors sold; it holds concentrated barrier
and dividend risk. It hurts when the index falls towards the barriers together, when gamma flips and hedging becomes
expensive (Book~5, chapter~27).
\iqlookfor{concentration of barrier risk and the desk's side of the trade.}
\end{solution}
